\titledquestion{Hash Table Operations}

Use \(E\) for an empty bin and \(D\) for a bin marked as deleted. Parts (b) and (c) concern the same table \(T\). The other parts are independent.

For open-addressing operations involving deleted bins, a search continues past a deleted bin and stops at the first empty bin or when the key is found. When inserting a key that is not already present, use the first deleted bin in its probing sequence if one occurs before the first empty bin; otherwise, use the first empty bin.

\begin{parts}

\part[2]
Suppose we wish to search a linked list of length \(n\), where each element contains a key \(k\) together with a hash value \(h(k)\). Each key is a long character string. How might we take advantage of the stored hash values when searching for a given key \(x\)?

\begin{solution}

When given a key $x$, we compute $h(x)$ firstly.

For every element in the linked list, we compare $h(x)$ with the stored hash value $h(k)$. If they are not equal, we skip to the next element. If they are equal, we then perform a full string comparison to confirm they are same character by character.

\end{solution}

\part[2]
Table \(T\) has capacity 11 and uses quadratic probing:
\[
h(k)=(5k+7)\bmod 11,\qquad
H_i(k)=(h(k)+i^2)\bmod 11,\qquad i=0,1,2,\ldots.
\]
Its current contents are:
\medskip
\begin{center}
\begin{tabular}{|c|c|c|c|c|c|c|c|c|c|c|c|}
\hline
Index&0&1&2&3&4&5&6&7&8&9&10\\\hline
Key&126&\(E\)&\(E\)&148&\(E\)&147&\(E\)&\(E\)&170&\(E\)&104\\\hline
\end{tabular}
\end{center}
\medskip

Execute \texttt{Insert 192}. Write the probing sequence (indices) and the resulting table.

\begin{solution}
$$
h(192)=10
$$

$$
H_0 = (10 + 0) \bmod 11 = 10
$$

$$
H_1 = (10 + 1^2) \bmod 11 = 11 \bmod 11 = 0
$$

$$
H_2 = (10 + 2^2) \bmod 11 = 14 \bmod 11 = 3
$$

$$
H_3 = (10 + 3^2) \bmod 11 = 19 \bmod 11 = 8
$$

$$
H_4 = (10 + 4^2) \bmod 11 = 26 \bmod 11 = 4
$$

Index 4 is empty, so we insert 192 there.

So the probing sequence is: 10, 0, 3, 8, 4

And the resulting table is:

\begin{center}
\begin{tabular}{|c|c|c|c|c|c|c|c|c|c|c|c|}
\hline
Index&0&1&2&3&4&5&6&7&8&9&10\\\hline
Key&126&\(E\)&\(E\)&148&192&147&\(E\)&\(E\)&170&\(E\)&104\\\hline
\end{tabular}
\end{center}

\end{solution}

\part[3]
Continue from the table obtained in part (b). Execute \texttt{Delete 148} using lazy deletion. Then perform \texttt{Search 192} and \texttt{Search 214}. For each search, write the probing sequence (indices) and whether the search succeeds. Finally execute \texttt{Insert 214} and write the resulting table.

\begin{solution}
After deleting 148, the table becomes:

\begin{center}
\begin{tabular}{|c|c|c|c|c|c|c|c|c|c|c|c|}
\hline
Index&0&1&2&3&4&5&6&7&8&9&10\\\hline
Key&126&\(E\)&\(E\)&\(D\)&192&147&\(E\)&\(E\)&170&\(E\)&104\\\hline
\end{tabular}
\end{center}

$$h(192) = 10$$

The probing sequence for searching 192 is: 10, 0, 3, 8, 4. The worth-noticing thing is we encounter $D$ at index 3, and the search continues. And after searching to index 4, we find 192, so the search succeeds.

$$h(214) = 10$$

$$
H_5(214) = (10 + 5^2) \bmod 11 = 35 \bmod 11 = 2
$$

The probing sequence for searching 214 is: 10, 0, 3, 8, 4, 2. When we reach index 2, we find an empty bin, so the search fails.

Then we insert 214. The probing sequence is: 10, 0, 3. When we reach index 3, we find $D$, so we insert 214 there.

The resulting table is:

\begin{center}
\begin{tabular}{|c|c|c|c|c|c|c|c|c|c|c|c|}
\hline
Index&0&1&2&3&4&5&6&7&8&9&10\\\hline
Key&126&\(E\)&\(E\)&214&192&147&\(E\)&\(E\)&170&\(E\)&104\\\hline
\end{tabular}
\end{center}

\end{solution}

\part[2]
Insert the letters
\[
C,\ O,\ M,\ P,\ U,\ T,\ I,\ N,\ G
\]
in that order into an empty hash table of capacity 12 using linear probing. The hash function is
\[
h(k)=5k\bmod 12,
\]
where \(k\) is the position of the letter in the alphabet: \(A=1,B=2,\ldots,Z=26\). Write the final table.

\begin{solution}

C($k = 3$): $h(3) = 5 \times 3 \bmod 12 = 15 \bmod 12 = 3$

O($k = 15$): $h(15) = 5 \times 15 \bmod 12 = 75 \bmod 12 = 3$. Collision with 3, store to index 4.

M($k = 13$): $h(13) = 5 \times 13 \bmod 12 = 65 \bmod 12 = 5$

P($k = 16$): $h(16) = 5 \times 16 \bmod 12 = 80 \bmod 12 = 8$

U($k = 21$): $h(21) = 5 \times 21 \bmod 12 = 105 \bmod 12 = 9$

T($k = 20$): $h(20) = 5 \times 20 \bmod 12 = 100 \bmod 12 = 4$. Collision with 4 and 5, store to index 6.

I($k = 9$): $h(9) = 5 \times 9 \bmod 12 = 45 \bmod 12 = 9$. Collision with 9, store to index 10.

N($k = 14$): $h(14) = 5 \times 14 \bmod 12 = 70 \bmod 12 = 10$. Collision with 10, store to index 11.

G($k = 7$): $h(7) = 5 \times 7 \bmod 12 = 35 \bmod 12 = 11$. Collision with 11, store to index 0.

The resulting table is:

\begin{center}
\begin{tabular}{|c|c|c|c|c|c|c|c|c|c|c|c|c|}
\hline
Index&0&1&2&3&4&5&6&7&8&9&10&11\\\hline
Key&G&\(E\)&\(E\)&C&O&M&T&\(E\)&P&U&I&N\\\hline
\end{tabular}
\end{center}

\end{solution}

\part[3]
For the keys
\[
91,\ 54,\ 67,\ 35,\ 69,\ 65,
\]
consider a hash table of capacity 6 with
\[
h_c(k)=((7k+3)\bmod c)\bmod 6.
\]
Find the smallest positive integer \(c\) for which the six keys have distinct hash values. Justify that your value is the smallest.

\begin{solution}

Let $y_k = 7k+3$.

$$
y_{91} = 7 \times 91 + 3 = 640
$$

$$
y_{54} = 7 \times 54 + 3 = 381
$$

$$
y_{67} = 7 \times 67 + 3 = 472
$$

$$
y_{35} = 7 \times 35 + 3 = 248
$$

$$
y_{69} = 7 \times 69 + 3 = 486
$$

$$
y_{65} = 7 \times 65 + 3 = 458
$$

They should not collide in a hash table of capacity 6, so $c \ge 6$.

When $c = 6$, we have $y_{35} = 248 \equiv 2 \pmod 6$ and $y_{65} = 458 \equiv 2 \pmod 6$, so $c \neq 6$.

When $c = 7$, $(7k + 3) \equiv 3 \pmod 7$, so $c \neq 7$.

When $c = 8$, we have $y_{35} = 248 \equiv 0 \pmod 8$ and $y_{69} = 486 \equiv 6 \pmod 8$, so $c \neq 8$.

When $c = 9$, we have

$$y_{91} = 640 \equiv 1 \pmod 9 \implies 1 \bmod 6 = 1$$

$$y_{54} = 381 \equiv 3 \pmod 9 \implies 3 \bmod 6 = 3$$

$$y_{67} = 472 \equiv 4 \pmod 9 \implies 4 \bmod 6 = 4$$

$$y_{35} = 248 \equiv 5 \pmod 9 \implies 5 \bmod 6 = 5$$

$$y_{69} = 486 \equiv 0 \pmod 9 \implies 0 \bmod 6 = 0$$

$$y_{65} = 458 \equiv 8 \pmod 9 \implies 8 \bmod 6 = 2$$

So $c = 9$ is the smallest positive integer for which the six keys have distinct hash values.

\end{solution}

\end{parts}
